\note{4}{Sat 25 May 2024 18:22}{Spectral Theory of Compact Operators}

This note will supplement more material on spectral theory. We will
\begin{itemize}
	\item Classify the spectrum and see how they interact with other operations
	\item Give concrete example to all concepts.
	\item Develop spectral theory for compact operators
\end{itemize}

\subsection{Spectral Theory}

\subsubsection{Classification}

\begin{definition}[Point spectrum, continuous spectrum, residual spectrum]
	\label{defn:PointSpectrumCont}
	Let $T: E \to E$ be an operator in Banach space. Let $\lambda \in \mathbb{C}$.
	\begin{enumerate}
		\item If $\operatorname{ker} (T - \lambda)$ contains a nonzero element, then $\lambda$ is an \textit{eigenvalue} of $T$, and we say $\lambda$ is in the \textit{point spectrum} of $T$, $\sigma_p(T)$.
		\item If $\operatorname{ker} (T - \lambda) = 0$ but $(T - \lambda)$ fails to be surjective, we say
			\begin{enumerate}
				\item $\lambda$ is in the \textit{continuous spectrum} of $T$, denoted $\sigma_c(T)$, if $\operatorname{Im} (T - \lambda)$ is dense.
				\item otherwise, we say $\lambda$ is in the \textit{residual spectrum} of $T$, denoted $\sigma_r(T)$.
			\end{enumerate}
	\end{enumerate}
	
\end{definition}

\begin{definition}[Essential Spectrum]
	\label{defn:EssentialSpectrum}
	A point $\lambda \in \mathbb{C}$ is in the \textit{essential spectrum} of $T$, $\sigma_{e}(T)$, if $T - \lambda$ fails to be Fredholm.
\end{definition}

\begin{proposition}
	\label{prop:TopEquivSpectrum}
	Topologically equivalent operators have the same $\sigma_p, \sigma_c, \sigma_r, \sigma_e$.
\end{proposition}

\begin{proposition}
	\label{prop:EssSpecCalkin}
	Let $H$ be a Hilbert space and $T \in B(H)$. Then the essential spectrum of $T$ is precisely the spectrum of the coset $T + \mathcal{K}(H)$ in the Calkin algebra.
\end{proposition}
\todo{proof omitted.}

\begin{proposition}
	\label{prop:ContEss}
	$\sigma_c(T) \subset \sigma_e(T)$
\end{proposition}
\begin{proof}
	Take $\lambda \in \sigma_c(T)$, then $\operatorname{Im} (T - \lambda)$ is dense in $E$ and $\operatorname{ker} (T - \lambda) = 0$. If $T - \lambda$ is Fredholm (i.e. $\lambda \not\in \sigma_e(T)$), then $\operatorname{Im} (T - \lambda)$ is closed, thus $T - \lambda$ surjective, hence a topological isomorphism. But this implies $\lambda \not\in \sigma(T)$, a contradiction.
\end{proof}


\begin{proposition}
	\label{prop:TopInjResidual}
	If there is $c>0$ such that $\|(T - \lambda) x\| > c \|x\|$ for all $x \in E$, that is, $T - \lambda$ is topologically injective, then $\lambda \in \sigma_r(T)$.
\end{proposition}
\todo{proof omitted.}

\subsubsection{Spectrum and Banach Dual}



\begin{proposition}[Spectrum of Banach Dual]
	\label{prop:SpectrumBanachDual}
	In $\mathbf{Ban}$, the spectrum of $T: E \to E$ coincides with the spectrum of the adjoint $T^*: E^* \to E^*$, and in addition,
	\begin{enumerate}
		\item If $\lambda \in \sigma_r(T)$, then $\lambda \in \sigma_p(T^*)$ ;
		\item If $\lambda \in \sigma_p(T)$, then $\sigma \in \sigma_p(T^*)$ or $\lambda \in \sigma_r(T)$ ;
		\item If $\lambda \in \sigma_c(T)$, then either $\lambda \in \sigma_c(T^*)$, or $\lambda \in \sigma_r(T^*)$. If $E$ reflexive, then $\lambda \in \sigma_c(T^*)$.
	\end{enumerate}
\end{proposition}
\begin{proof}
	We first show $\sigma(T) = \sigma(T^*)$. Since the adjoint functor preserves isomorphism, we have $T - \lambda$ invertible $\implies T^* - \lambda$ invertible $\implies T^{* *} - \lambda$ invertible. By considering the canonical embedding $E \hookrightarrow E^{* *}$ we have $T^{* *} - \lambda$ invertible $\implies T -  \lambda$ invertible. (A side note: $(T - \lambda)^* = T^* - \lambda$ since we are considering the Banach adjoint, not the Hilbert one.)

	(1). Take $\lambda \in \sigma_r(T)$, then $\operatorname{Im} (T - \lambda)$ is not dense. Take $x \not\in \overline{\operatorname{Im} (T - \lambda)} $, and let $f \in E^*$ be such that $f(x) \neq 0$ and $f|_{\operatorname{Im} (T - \lambda)} = 0$. Now $(T - \lambda)^* f = f \circ (T - \lambda) = 0$, thus $\operatorname{ker} (T - \lambda)^* \neq 0$, i.e. $\lambda \in \sigma_p(T^*)$.

	(2) Take $x \in \operatorname{ker} (T - \lambda)$, $x \neq 0$. We show $\lambda \not\in \sigma_c(T^*)$. For all $f \in E^*$ we have $ ((T - \lambda)^* f)x = f\left( (T - \lambda) x \right)  =0$. Thus for $g \in E^*, g(x) \neq 0$ we have $g \not\in \overline{\operatorname{Im} (T^* - \lambda)} $.

	(3) Take any $f \in \operatorname{ker} (T - \lambda)^*$. Then $(T^* - \lambda) f = f\circ (T - \lambda) = 0$. But $\operatorname{Im} (T - \lambda)$ is dense, so $f = 0$, i.e. $\lambda \not\in \sigma_p(T^*)$. If $E$ is reflexive, then $\lambda \in \sigma_r(T^*) \implies \lambda \in \sigma_p(T^{* *}) = \sigma_p(T)$ which is a contradiction.
\end{proof}


\subsubsection{Spectrum and Hilbert Space Dual}

\begin{proposition}[Spectrum of Hilbert Dual]
	\label{prop:SpectrumHilbertDual}
	Consider $T: H \to H$ in $\mathbf{Hilb}$. We have $\sigma(T^*) = \{\overline{\lambda} : \lambda \in \sigma(T)\} $, and
	\begin{enumerate}
		\item If $\lambda \in \sigma_r(T)$, then $\overline{\lambda}  \in \sigma_p(T^*)$;
		\item If $\lambda \in \sigma_p(T)$, then $\overline{\lambda}  \in \sigma_p(T^*) \cup \sigma_r(T^*)$;
		\item $\sigma_c(T^*) = \{\overline{\lambda} : \lambda \in \sigma_c(T)\} $ ;
		\item $\sigma_e(T^*) = \{\overline{\lambda} : \lambda \in \sigma_e(T)\} $.
	\end{enumerate}
\end{proposition}
\begin{proof}
	The results follow from Proposition~\ref{prop:SpectrumBanachDual} and the fact that Hilbert space is reflexive. Use the commutative diagram relating Hilbert adjoint to Banach adjoint.

	The proof can also be made directly by repeating the argument in Banach adjoint.
\end{proof}


\begin{proposition}
	\label{prop:NormalResidualSpectrum}
	Normal operators have no residual spectrum.
\end{proposition}


\subsubsection{Spectral Theorem for Compact Operator}


\begin{theorem}[Riesz]
	\label{thm:RieszCompact}
	Let $T: E \to E$ be a compact operator acting in an infinite dimensional space. Then 
	\begin{itemize}
		\item $0 \in \sigma(T)$, 
\item $\sigma_e(T) = \{0\} $.
\item For each $\lambda \in \sigma(T) \setminus \{0\} $, $\operatorname{dim} \operatorname{ker} (T - \lambda) < \infty $.
\item $\sigma(T)$ at most countable, and $0$ is the only possible limit point.
	\end{itemize}
\end{theorem}
\begin{proof}
	Since $T$ is compact, $T$ is not Fredholm (Proposition~\ref{prop:CompactNotFred}), so $0 \in \sigma_e(T)$. By the Fredholm Alternative, for each $\lambda \in \sigma(T), \lambda\neq 0$, $T - \lambda$ is Fredholm and $\operatorname{dim} \operatorname{ker} (T - \lambda) < \infty $. Thus $\sigma_e(T) = \{0\} $.

	Now we show $\sigma(T)$ is at most countable, and $0$ is its only possible limit point. If either were false, then there exists a sequence $\lambda_n$ of points in the spectrum that tends to some $\theta\neq 0$. Without loss of generality we may assume $|\lambda_n| > \frac{|\theta|}{2}$, and all $\lambda_n$ are distinct. We know all the $\lambda_n$ are eigenvalues. Then we may take a unit vector in each eigenspace, $x_n \in \operatorname{ker} (T - \lambda_n)$. From linear algebra, each $x_n$ are linearly independent. Set $E_n = \operatorname{span } \{x_1, \ldots, x_n\} $.

	Now use Riesz lemma of perpendicularity (Lemma~\ref{lem:Riesz}), we choose for every $n > 1$, $y_n \in E_n$ such that $\|y \| = 1$ and $\text{dist}(y_n, E_{n - 1}) > \frac{1}{2}$. Write $y_n = \mu_n x_n + z_n$ where $\mu_n \in \mathbb{C}, z_n \in E_{n - 1}$. Then for $n > m$, we have $\|T y_n - T y_m\| =\|T (\mu_n x_n + z_{n - 1}) - T y_m\| = \|\lambda (\mu_n x_n) - T(y_m - z_{n - 1})\| > |\lambda| \text{dist}(y_n, E_{n - 1}) > \frac{|\lambda|}{2} > \frac{|\theta|}{4}$. This contradicts compactness of $T$.
\end{proof}

\begin{corollary}[Compact Operator in Hilbert Space]
	\label{cor:CompactHilbert}

	Let $H$ be a Hilbert space and let $T \in \mathcal{K}(H)$. Then $|T| \in \mathcal{K}(H)$. If the spectrum of $|T|$ is listed as $\lambda_1 \ge \lambda_2 \ge  \ldots$ where each eigenvalue is repeated according to its multiplicity, then there are orthonormal families $(x_n), (y_n)$ consisting of vectors in $H$ such that
	\[
		T = \sum_{k\ge 1} \lambda_k \left\langle \cdot , x_k \right\rangle y_k
	\] 
	where the series is convergent in the norm of $B(H)$.
\end{corollary}
\begin{proof}
	We show $|T|$ is compact. Of course $|T|^2 = T T^*$ is compact. Now both of them are positive and are diagonalizable; the claim follows from Proposition~\ref{prop:DiagonalOperatorCompact}. \ask{check.}
	By the above theorem, $\sigma(|T|)$ is countable. Use the spectral theorem for normal operators (Proposition~\ref{prop:SpectrumNormalOperator}) to write $|T|$ as
	\[
		T 
		= \sum_{n\ge 1} \lambda_n E(\lambda_n)
		= \sum_{k\ge 1} \lambda_k \left\langle \cdot , x_k \right\rangle x_k
	.\] 
	Now $\sigma(|T|) \subset [0,\infty )$ and $0$ is the only possible accumulation point; thus we can write $\lambda_k$ in descending order.

	Now apply the polar decomposition (Theorem~\ref{thm:PolarDecomposition}) to $T$ and have $T = V |T|$ ; the partial isometry $V$ has the form (Proposition~\ref{prop:MatrixPartialIsometry}) $V = \sum_{n \geq 1}\left\langle\cdot, x_n\right\rangle y_n$. Thus $T=\sum_{k \geq 1} \lambda_k\left\langle\cdot, x_k\right\rangle y_k$.
\end{proof}

There are several other results of the same type, where one is interested in decomposing a compact operator in a Hilbert space. Those results will be immediate corollaries since we have now in hand the spectral theory of normal operators in Hilbert space; but they can also be developed without spectral theory as done in Khelemskii 2006. Here we will omit most of the proofs.

\begin{theorem}[Schmidt]
	\label{thm:Schmidt}
	Let $H$ and $K$ be Hilbert spaces and $T: H \to K$ be a compact operator. Then there exist
	\begin{enumerate}
		\item an orthonormal system $(e_i')$ in $H$ of finite or countable cardinality;
		\item  an orthonormal system $(e_i'')$ in $K$ of the same cardinality and
		\item a (finite or infinite) sequence $s_1 \ge s_2 \ge \ldots$ of positive numbers with the index set of same cardinality, tending to zero if it is infinite,
	\end{enumerate}
	such that $T$ can be represented in the form
	\[
		T = \sum_n s_n \left\langle \cdot , e'_n \right\rangle e''_n
	.\] 
\end{theorem}

\begin{theorem}[Hilbert-Schmidt]
	\label{thm:HilbertSchmidt}
	Let $T: H \to H$ be compact self-adjoint operator on a Hilbert space. Then there exist
	\begin{enumerate}
		\item an orthonormal system $(e_i')$ in $H$ of finite or countable cardinality, and
		\item a (finite or infinite) sequence $\lambda_1 , \lambda_2 , \ldots$ of nonzero real numbers with the index set of same cardinality, tending to zero if it is infinite,
	\end{enumerate}
	such that $T$ can be represented in the form
	\[
		T = \sum_n \lambda_n \left\langle \cdot , e_n \right\rangle e_n
	.\] 
\end{theorem}

The following is an equivalent formulation of the Hilbert-Schmidt theorem.

\begin{proposition}
	\label{prop:HilbertSchmidtEquiv}
	Let $T: H \to H$ be compact self-adjoint operator on a Hilbert space. Then there exist a finite or infinite sequence $\lambda = (\lambda_1, \lambda_2, \ldots)$ of real numbers, tending to zero if it is infinite, and a Hilbert space $H_0$ such that $T$ is unitarily equivalent to the operator $R: l_2^m \oplus H_0 \to  l_2^m \oplus H_0$, $m$ being cardinality of $\lambda$, that acts on $l_2^m$ as the diagonal operator $T_\lambda$ and take $H_0$ to zero. In short,
	\[
		T = T_\lambda \oplus 0
	.\] 
\end{proposition}

For separable Hilbert space the Hilbert-Schmidt theorem can be simplified to

\begin{proposition}
	\label{prop:HilbertSchmidtSep}
	Let $T: H \to H$ be compact self-adjoint operator in a separable Hilbert space. Then
	\begin{enumerate}
		\item there is an orthonormal basis in $H$ consisting of eigenvectors of $T$. The corresponding sequence of eigenvalues is real and tends to zero;
		\item $T$ is unitarily equivalent to the diagonal operator $T_\lambda: l_2^m \to l_2^m$, where $m$ is Hilbert dimension of $H$, and $\lambda$ is the sequence of eigenvalues as in (1).
	\end{enumerate}
\end{proposition}


\subsubsection{Examples in Lp and lp}

Now we compute the spectrum of a number of concrete operators.

\begin{proposition}
	Let $T_\mu; \mu = (\mu_1, \mu_2, \ldots)$ be the diagonal operator on $l_p; 1 \le p \le \infty $. Then
	\begin{itemize}
		\item $\sigma(T) = \overline{\{\mu_n\} } $,
		\item $\sigma_p(T) = \{\mu_n\} $,
		\item $\sigma(T) \setminus \sigma_p(T) = \sigma_c(T)$ if $p < \infty $,
		\item $\sigma(T) \setminus \sigma_p(T) = \sigma_r(T)$ if $p = \infty $,
		\item $\sigma_e(T)$ is set of limit points of $\mu$.
	\end{itemize}
\end{proposition}
\begin{proof}
	
	The statement about $\sigma(T)$ and $\sigma_p(T)$ are clear.
	We prove the statements regarding $\sigma_c(T)$ and $\sigma_r(T)$. Take $S = T - \lambda, \lambda \in \sigma(T) \setminus \sigma_p(T)$. When $p < \infty$, $S$ is nonzero coordinate-wise, and for any element $x \in l_p$ there exists a finite-dimensional truncation $x'$ close to $x$, where $x' \in \operatorname{Im} S$.  Thus $\operatorname{Im} S$ is dense in $l_p$. For $p = \infty $, for any $x \in l_\infty $, $S x$ gets close to $0$ i.o., thus $\text{dist}(\xi, \operatorname{Im} S) \ge 1$ where $\xi= (1,1, \ldots)$. Thus  $\operatorname{Im} S$ not dense in $l_\infty $.
\end{proof}


\begin{proposition}
	Let $T_f: L_p(X, \mu) \to L_p(X, \mu); f \in L_\infty (X, \mu)$ be multiplication operator by $f$. Then
	\begin{itemize}
		\item $\sigma(T_f) = f(X) _{\text{ess}}$ essential values of $f$ (numbers $\lambda$ such that for every $\varepsilon>0$, $\mu \{|f(t) - \lambda| < \varepsilon\}  > 0$);
		\item $\sigma_p(T_f) = \{\lambda: \mu \{f(t) = \lambda\} > 0\} = \{\lambda: \mu(Z_\lambda) > 0\} $;
		\item $\sigma(T_f) \setminus \sigma_p(T_f) = \sigma_c(T_f)$ if $p < \infty $;
		\item $\sigma(T_f) \setminus \sigma_p(T_f) = \sigma_r(T_f)$ if $p = \infty $ ;
		\item If $(X, \mu)$ is an interval $[a,b]$ with Lebesgue measure, $\sigma_e(T_f) = \sigma(T_f)$.
	\end{itemize}
\end{proposition}
\begin{proof}
	
	Recall the bounded multiplication operators are exactly those $T_f$ where $f \in L_\infty (X, \mu)$. If $\lambda \in f(X)_{\text{ess}}$ then $f^{-1}$ DNE or unbounded, thus $\lambda \in \sigma(T_f)$. On the other hand, if $\lambda \not\in f(X)_{\text{ess}}$ then $T_{f^{-1}}$ is bounded and $\lambda \not\in \sigma(T_f)$.

	$T_f - \lambda$ is injective iff $\mu(Z_\lambda) = 0$. Thus $\sigma_p(T_f) = \{\lambda: \mu(Z_\lambda) > 0\} $.

	Take $\lambda \in \sigma(T_f) \setminus \sigma_p(T_f)$. Suppose $p < \infty $, then $T_f - \lambda$ is injective, and  every $g$ with support on $\{|f - \lambda| \ge \varepsilon\} $ is mapped onto. Thus $\operatorname{Im} (T_f - \lambda)$ is dense, whence $\lambda \in \sigma_c(T_f)$. When $p = \infty $, we have $\|(T_f - \lambda) \xi - 1_X\| \ge 1$, where $1_X$ is constant function $1$ on $X$ and $\xi \in L_\infty(X, \mu)$. Thus $1_X$ is separated from $\operatorname{Im} (T_f - \lambda)$, whence $\lambda \in \sigma_r(T_f)$.

	Since $(T_f - \lambda) x = T_{f - \lambda}x$, $\sigma_e(T_f) = \sigma(T_f)$ follows from Proposition~\ref{prop:LpMultFredholm}.
\end{proof}


\begin{proposition}
	Let $T_l, T_r: l_p \to l_p; 1 \le p \le \infty $ be the left and right shift operators. We have
	\begin{itemize}
		\item $\sigma_p(T_l) = \begin{cases}
				\mathbb{D}^0 & p < \infty \\
				\mathbb{D} & p = \infty 
		\end{cases}$ 
		\item $\sigma_p(T_r) = \emptyset $
		\item $\sigma_r(T_r) = \begin{cases}
				\mathbb{D}^0 & 1 < p < \infty \\
				\mathbb{D} & p = 1 \\
				\text{contains } \mathbb{D}^0 & p = \infty 
		\end{cases}$ 
	\end{itemize}
	Note: when $p = \infty $ we have $\sigma_r(T_r) = \mathbb{D}$, we will show this after the next lemma.

\end{proposition}
\begin{proof}
	
	Let's first compute $\sigma_p(T_l)$. $T_\lambda \xi = \lambda \xi$ iff $\xi = c\left( 1, \lambda, \lambda^2, \ldots \right) $, which is in $l_p$ iff $\lambda \in \mathbb{D}^0$  when $p < \infty $, and iff $\lambda \in \mathbb{D}$ when $p = \infty $.

	$\operatorname{Im} T_r$ has a nontrivial kernel, which consists of all functions with 0 in the first coordinate.

	Finally, use Proposition~\ref{prop:SpectrumBanachDual} and the fact that $(l_p)^* = l_q$ where $\frac{1}{p} + \frac{1}{q} = 1$, to arrive at result about $T_r = T_l^*$.
\end{proof}




\begin{lemma}
	\label{lem:ShiftTopEquiv}
	Let $T$ be any of the shift operators on $l_p, l_p(\mathbb{Z}),$ or in $L_p(\mathbb{R})$ (in this case $T = T_a, a \neq 0$ ), assuming that  $1 \le p \le \infty $. Then for each $\lambda \in \mathbb{T}$, the operators $T - 1$ and $T - \lambda$ are weakly topologically equivalent. As a corollary, each of the sets $\sigma(T), \sigma_p(T), \sigma_c(T), \sigma_r(T)$, and $\sigma_e(T)$ either contains the entire $\mathbb{T}$ or is disjoint from $\mathbb{T}$.
\end{lemma}
\begin{proof}
	In the case $l_p$, take $T_l$ to be the left shift operator, and let $\alpha = (1, \lambda, \lambda^2, \ldots)$ and $\beta = (\lambda, \lambda^2, \lambda^3, \ldots)$. Let $T_\alpha, T_\beta$ be diagonal operators corresponding to $\alpha, \beta$. The following diagram
% https://quiver.local:8080/#q=WzAsNCxbMCwwLCJcXGVsbF9wIl0sWzIsMCwiXFxlbGxfcCJdLFsyLDIsIlxcZWxsX3AiXSxbMCwyLCJcXGVsbF9wIl0sWzAsMSwiVF9sIC0gMSJdLFszLDIsIlRfbCAtIFxcbGFtYmRhIl0sWzAsMywiVF9cXGFscGhhIiwxXSxbMSwyLCJUX1xcYmV0YSIsMV1d
\[\begin{tikzcd}[ampersand replacement=\&]
	{\ell_p} \&\& {\ell_p} \\
	\\
	{\ell_p} \&\& {\ell_p}
	\arrow["{T_l - 1}", from=1-1, to=1-3]
	\arrow["{T_l - \lambda}", from=3-1, to=3-3]
	\arrow["{T_\alpha}"{description}, from=1-1, to=3-1]
	\arrow["{T_\beta}"{description}, from=1-3, to=3-3]
\end{tikzcd}\]
commutes.
Similar construction holds for all shift operators in $l_p$ and $ l_p(\mathbb{Z})$. For  $L_p(\mathbb{R})$, we are inspired by the $l_p$ case to postulate that the diagram
% https://quiver.local:8080/#q=WzAsNCxbMCwwLCJMX3AoXFxSKSJdLFsyLDAsIkxfcChcXFIpIl0sWzAsMiwiTF9wKFxcUikiXSxbMiwyLCJMX3AoXFxSKSJdLFswLDEsIlRfYSAtIDEiXSxbMiwzLCJUX2EgLVxcbGFtYmRhIl0sWzAsMiwiZV57aSBjXzEgdH0iLDFdLFsxLDMsImVee2kgY18yIHQgKyBpIFxcdGhldGF9IiwxXV0=
\[\begin{tikzcd}[ampersand replacement=\&]
	{L_p(\R)} \&\& {L_p(\R)} \\
	\\
	{L_p(\R)} \&\& {L_p(\R)}
	\arrow["{T_a - 1}", from=1-1, to=1-3]
	\arrow["{T_a -\lambda}", from=3-1, to=3-3]
	\arrow["{e^{i c_1 t}}"{description}, from=1-1, to=3-1]
	\arrow["{e^{i c_2 t + i \theta}}"{description}, from=1-3, to=3-3]
\end{tikzcd}\]
commutes. Taking $f(t) \in L_p(\mathbb{R})$ and doing some algebra yields a unique solution $\theta = \text{arg} \lambda, c_1 = c_2 = \frac{\theta}{a}$.
\end{proof}



\begin{proposition}
	\begin{itemize}
		\item In $l_\infty $ we have $\sigma_r(T_r) = \mathbb{D}$.
	\end{itemize}

\end{proposition}
\begin{proof}
	
	We already know $\sigma_r(T_r) \supset \mathbb{D}^0$. We only need to show $\sigma_r(T_r) \supset \mathbb{T}$. In view of Lemma~\ref{lem:ShiftTopEquiv}, we only need to show $T_r - 1$ has no dense image (it's injective since $\sigma_p(T_r) = \emptyset $ ). Take any $\eta = (T_r - 1) x \in \operatorname{Im} (T_r - 1)$, and we have $\lim_{n \to \infty} \frac{1}{n} \sum_{k=1}^{n} \eta_k = 0$. This implies $\|\eta - (1, 1, \ldots)\|_\infty \ge 1$, that is, $\operatorname{Im} (T_r - 1)$ is not dense.
\end{proof}

\begin{proposition}
	\label{prop:SpectShift}
	The spectrum of left or right shift operators on $l_p; 1 \le p \le \infty $ is $\mathbb{D}$, and\\
	(i) in the case of $l_p ; 1<p<\infty$ and $c_0$ we have $\sigma_p\left(T_l\right)=\sigma_r\left(T_r\right)=\mathbb{D}^0$, $\sigma_c\left(T_l\right)=\sigma_c\left(T_r\right)=\mathbb{T}$, and $\sigma_r\left(T_l\right)=\sigma_r\left(T_r\right)=\varnothing$;\\
(ii) in the case of $l_1$ we have $\sigma_p\left(T_l\right)=\mathbb{D}^0, \sigma_c\left(T_l\right)=\mathbb{T}, \sigma_r\left(T_r\right)=\mathbb{D}$, and $\sigma_r\left(T_l\right)=\sigma_p\left(T_r\right)=\sigma_c\left(T_r\right)=\varnothing ;$\\
(iii) in the case of $l_{\infty}$ we have $\sigma_p\left(T_l\right)=\sigma_r\left(T_r\right)=\mathbb{D}$, and $\sigma_c\left(T_l\right)=$ $\sigma_c\left(T_r\right)=\sigma_r\left(T_l\right)=\sigma_p\left(T_r\right)=\varnothing$.
\end{proposition}
\begin{proof}
	By Proposition~\ref{prop:SpectrumClosed}, $r(T) \le \|T\| = 1$, and the spectrum is closed containing $\mathbb{D}^0$, hence equals to $\mathbb{D}$. Now, in the case $1 \le p < \infty $, points in $\mathbb{T}$ are eigenvalue of $T_l$ or $T_r$, so by Proposition~\ref{prop:SpectrumHilbertDual} it cannot be in the residual spectrum, i.e. it must be in the continuous spectrum.
\end{proof}





\begin{claim}
	Let $T_b$ be bilateral shift in $l_p(\mathbb{Z})$ or $L_p(\mathbb{R})$.
	\begin{itemize}
		\item When $p = 1, \sigma_r(T_b) \supset \mathbb{T}$,
		\item When $p = \infty , \sigma_p(T_b) \supset \mathbb{T}$.
	\end{itemize}
\end{claim}
\begin{proof}
	
	Using duality of $l_p$ and $L_p$ spaces, the statement about $p = 1$ implies the statement about $p = \infty $. Consider $l_1(\mathbb{Z})$. The composition of the following to maps is zero:
	\[
		l_1 \xrightarrow{T - 1} l_1 \xrightarrow{f} \mathbb{C}
	\] 
	where $f: \eta \mapsto \sum \eta_n$. Thus $\operatorname{Im} (T - 1) \subset \operatorname{ker} f$ which is not dense in $l_1(\mathbb{Z})$.

	For $L_1(\mathbb{R})$, use $f: \xi \mapsto \int_{\mathbb{Z}} \xi(t) d t$.
\end{proof}

\begin{claim}
	The point spectrum of the shift operator by $a \in \mathbb{T}$ in $L_p(\mathbb{T})$ contains all numbers $a^n; n \in \mathbb{Z}$.
\end{claim}
\begin{proof}
	Embed $\mathbb{T} \subset \mathbb{C}$ and identify $T_a$ as $a \in \mathbb{T}$ acting by complex multiplication. Then on the function $f(z) = z^n$ in $ L_p(\mathbb{T})$, we have $T_a f (z) = f(a z)= a^n f(z)$.
\end{proof}



\begin{proposition}
	\label{prop:SpectBiShift}
	For all $l_p; 1 \le p \le \infty $, the spectrum of bilateral shift $l_p:(\mathbb{Z}) \to l_p(\mathbb{Z})$ and $T_a:L_p(\mathbb{R}) \to L_p(\mathbb{R}); a \neq 0$ is $\mathbb{T}$. Furthermore,
	\begin{itemize}
		\item when $1 < p < \infty $, $\sigma(T_b) = \sigma_c(T_b) = \mathbb{T}$
		\item when $p=1 $, $\sigma(T_b) = \sigma_r(T_b) = \mathbb{T}$
		\item when $p=\infty  $, $\sigma(T_b) = \sigma_p(T_b) = \mathbb{T}$
	\end{itemize}
\end{proposition}
\begin{proof}
	The cases $p = 1, p = \infty $ have already been established.

	Consider the case $1 < p < \infty $. We can explicitly check that for $|\lambda|\neq 1$, $T - \lambda$ is invertible. By Lemma~\ref{lem:ShiftTopEquiv} and that $\sigma(T_b)$ is nonempty, $\sigma(T_b) = \mathbb{T}$. Now $1$ is eigenvalue of our operator iff $p = \infty $. Thus $\sigma_p(T_b) = \emptyset $ for all $p \in (1, \infty )$. But by duality property, $\sigma_r(T_b) = \emptyset $ for $p \in (1, \infty )$. Thus $\sigma(T_b) = \sigma_c(T_b) = \mathbb{T}$. 
\end{proof}

\begin{proposition}
	\label{prop:SpectTShift}
	Spectrum of shift operator by an element $a \in \mathbb{T}$ in $L_p(\mathbb{T})$ is as follows:
	\begin{itemize}
		\item if $a = e^{2 \pi i \frac{m}{n}}$ where $m, n$ coprime, then $\sigma(T_a) = \sigma_p(T_a) = \{e^{2 \pi i \frac{k}{n}}; k = 0, 1, \ldots, n - 1\} $.
		\item in the other cases $\sigma(T_a) = \mathbb{T}$.
	\end{itemize}
\end{proposition}
\begin{proof}
	In the first case, we have $(T_a)^n = 1$, so the Spectral Mapping Theorem implies $\sigma(T_a)$ must be the $n$-th roots of unity. But we have just shown that $\sigma_p(T_a)$ consists of $\{a^n; n \in \mathbb{Z}\} $, i.e. all the $n$-th roots of unity.

	In the second case, $\{a^n; n \in \mathbb{Z}\} $ is dense in $\mathbb{T}$, so the result follows from the fact that $\sigma(T_a)$ is closed.
\end{proof}

\begin{remark}
	(Khelemskii page 318) In the ``Hilbert'' case $p = 2$, every bilateral shift operator in $l_2(\mathbb{Z}), L_2(\mathbb{R}), L_2(\mathbb{T})$ is unitarily equivalent to some operator of multiplication by a function, established by the Fourier operators. We will give the examples later.
\end{remark}

\subsubsection{Topologically nilpotent operators}

\begin{definition}[Quasinilpotent Operators]
	\label{defn:QuasinilpotentOperators}
	We say an element $a$ of a Banach algebra is \textit{topologically nilpotent} or \textit{quasinilpotent} if $\lim_{n \to \infty} \|a^n\|^{\frac{1}{n}} = 0$.

\end{definition}

By the spectral radius formula, quasinilpotent elements in a Banach algebra are precisely those for which the spectrum consists of zero only.

\begin{proposition}[Volterra Operator]
	\label{prop:SpectVolterroperator}
	Let $T$ be operator of indefinite integration in $L_1[0,1], L_2[0,1]$ or $C[0,1]$. This is a topologically nilpotent operator, so $\sigma(T) = \{0\} $.

	The same is true for the Volterra operator in $L_2[a,b]$ with essentially bounded kernel $K(s,t)$. As a corollary, the integral equation
	\[
		\int _a^t K(s, \tau) x(\tau) d \tau - \lambda x(s) = y(s)
	\] 
	has a unique solution in $L_2[a,b] $ for each $\lambda \in \mathbb{C} \setminus \{0\} $ and for every right-hand side $y \in L_2[a,b]$.
\end{proposition}
\begin{proof}
	We show $\|T^n\|^{1 / n} \to 0$. In each space, if we take $f$ inside the unit ball, we would have $T f(x) \le \|f\| \le 1$. Apply $T$ to both sides, we have $T^{n+1} f(x) \le \frac{t^n}{n!}; t \in [0,1]$. Thus $\|T^n\| \le \frac{1}{n!}$.
\end{proof}



\begin{proposition}
	\label{prop:NormIndefInt}
	The norm of the operator of indefinite integration on $L_2[0,1]$ is $\frac{2}{\pi}$.
\end{proposition}
\begin{proof}
	We write $T = U S$, where
	\[
		S x(s) = \int _0^{1 - s} x(t) d t , \qquad
		U x(s) = x(1 - s)
	.\] 
	Since $U$ is unitary, $\|T\| = \|S\|$. We know $S$ is self-adjoint, so $\|S\| = r(S)$. To find eigenvalues of $S$ we search for non-zero solutions to the equation
	\[
		S x = \lambda x
	.\] 
	The solution to this equation is exactly solution to the ODE
	\[
		x(1 - t) + \lambda x'(t) = 0, \quad x(1) = 0
	.\] 
	Differentiate again and substitute back the first-order equation, we see that it suffices to find solutions to
	\[
		x(t) + \lambda^2 x''(t) = 0
	.\] 
	We understand this ODE very well; it has solution $x(t) = \cos \left( \left( \frac{\pi}{2} + k \pi \right) t \right)$, when $\lambda = \left( \frac{\pi}{2} + k \pi \right) ^{-1}$, for each $k \in \mathbb{Z}$.
	
	Thus $|\lambda| \le \frac{2}{\pi}$, and this maximum is attained since $S x = \frac{2}{\pi} x$ for $x(t) = \cos \left( \frac{\pi}{2} t \right) $.

\end{proof}

